\section{Постановка}
\label{sec:problem}

\porozhdeno{\Claim{P01}{Цель работы --- оценить, насколько таксономические назначения и профили одних и тех же почвенных образцов зависят от выбора готовой модели классификатора. Проводится описательное сравнение моделей, связанных с базами SILVA, GTDB и Greengenes2 (далее GG2), с использованием единой процедуры таксономического назначения \texttt{classify-sklearn} и фиксированных входных последовательностей. Оцениваются покрытие таксономическими назначениями, буквальное совпадение таксономических линий и расстояния между профилями. В отсутствие эталонных данных сравнение не позволяет определить наиболее точную модель.}}

\porozhdeno{\Claim{P02}{Набор данных включает \ProjectSamplesSoil{} почвенных образцов из московских районов Крылатское и Люблино и отрицательный контроль K\_S63 (вода). Контроль обработан совместно с почвенными образцами, но исключён из расчёта сводных показателей по почвам. Таблица, полученная после денойзинга, содержит \ProjectAsvAll{} вариантов ампликонных последовательностей (ASV), из которых \ProjectAsvSoil{} обнаружены в почвенных образцах. КР1 и ЛЮБ1 обозначают точки отбора: с КР1 связан Kr1\_S80, с ЛЮБ1 --- LUB2\_S92, LUB3\_S93 и LUB4\_S94; химические измерения относятся к отдельной пробе каждой точки. Единицей сравнения таксономических профилей служит образец, а единицей сравнения назначений --- ASV. Отдельные ASV и технические повторы не рассматриваются как независимые биологические повторности.}}

\porozhdeno{\Claim{P03}{Правило сравнения зафиксировано в методике 10 ноября 2026 года, до выполнения сравнительного расчёта. Постановка задачи и выбор результатов для представления документированы 18 ноября 2026 года. Последующий разбор результатов имеет описательный характер и не является проверкой заранее зарегистрированных гипотез.}}
